% !TeX root = ../papers/nc-p3.tex
\section{滑らかな判別式の全持上げの成分と固定計量例外}
\label{sec:marked-smooth-discriminant-components}
\begin{theorem}\label{thm:marked-smooth-discriminant-components}
$P=Q$、$G,F$ が対称で、判別三次曲線が滑らかな正規化印付き二切断表示の空間は
既約で次元27である。全ての持上げは分裂混合点を通る片側対称族の同じ閉包に入り、
基底変更を許した全窓は既知の十五次元成分に属する。一般点の窓は全次数正則・四辺有限・
反対階数10、最小周期4で射影自己同型群は自明、添字保持同型類の像は十四次元である。

滑らかな $E$ と計量を共に固定すれば、定理\ref{thm:marked-fixed-metric-elliptic}の
例外外では\emph{全て}の持上げの表示空間は既約で次元18である。
従って八次元の一般正則な同型類の族の閉包に、可換二平面を持つ印での解も全て含まれる。
固定計量例外では次元18の既約な正規化表示成分はちょうど四つである。
例外を与える平方根を $M_0$ とすれば、それらは $-2d=H-M_0$ の四つの位数4の点で区別される。
この最後の四成分の個数は座標と印を固定した二切断表示でのものであり、
各成分の一般AS性や、印を忘れたときの個数を前提としない。
\end{theorem}
\begin{proof}[相手の動きに沿う右核線束]
対称三平面 $G',F'$ の順で混合行列を作り、列側を $F'$ とする。
滑らかな行列式があれば両三平面は既約である。共通不変平面は自己随伴性で共通不変直線を与え、
その直線の代表 $n$ を保つ三平面は $n\odot n^\perp$ を消す跡零三平面そのものとなる。
その二次鉛筆は因子 $n\cdot v$ を持つ。この直線上の全 $T_v$ が三平面に入るので
混合行列式にも直線因子が出る。二空間の和の次元も5である。
前節までの $E$ 上の部分束 $\mathcal U$ を用いると
$\det\mathcal U^*=3H-B$ である。実際 $X(v)\wedge T_v$ の三次成分は、
独立交換子の場合には鉛筆基点のイデアルを生成する。
可換な場合も $X$ と $T_v$ が張る飽和部分束の計算を以下で直接行う。

独立交換子の場合、中心が有限なので節\ref{sec:marked-cubic-linear-system}の
安定性・五切断の証明は非被約 $B$ にもそのまま成立する。
$F'$ に対応する次数二線束 $M\subset\mathcal U^*$ の二切断は
$F'^\perp$ の制限である。そこで $\mathcal U\cap(F'\otimes\mathcal O_E)$ の直線は
$(\mathcal U^*/M)^*$ であり、一般点では混合行列の右核そのものである。
両者は飽和した線束なので、右核の射影写像が与える線束は
\[
             T_C=3H-B-M,\qquad d_C=2H-B-M.
\]
従って $F'$ の楕円五次曲線から $\mathrm{Pic}^0(E)$ への $d_C$ は
平行移動を伴う反転同型である。五基点を除いても、$-2$ の逆像の有限エタール被覆は連結である。

可換二平面の場合には $\mathcal U$ 内で $X(v)$ が生成する線を飽和すると
$-H+p$ となる。$X$ の二つの独立座標は $p$ を切り、滑らかな曲線上で共通零は長さ1である。
$T_v$ との外積の共通因子は曲線上でちょうど $B=T+p$ になることを確認する。
$p$ の一次零を除いた $X$ を $\bar X$ とする。$\bar X\wedge T_v$ は $T$ の外では消えない。
$T$ の各点で $T_v\bmod H_0$ の成分が生成するイデアルは $\mathcal I_T\mathcal O_E$ である
（$\ker R$ は $I_T(2)$ 全体）。よって外積の消失次数はそこでの $T$ の長さ以下である。
また $T$ 上で $T_v$ は $H_0$ にあり、$p$ の外では $X$ と比例する。
$p\in T$ の場合にも $B$ 上では $T_v$ が $G'$ にあり、$G'$ 上の交代三行列は $T_v$ を消す。
この条件は $X\wedge T_v=0$ と同値であり $B$ 上で成立する。
従って $p$ の一次零を引いた後も $T$ の全長で外積は消え、主張の共通因子を得る。
ゆえにこの場合も $\det\mathcal U^*=3H-B$ である。
$F'^\perp$ の二切断は飽和した $X$ への制限により $H-p$ の完全で基点自由な二切断に写る。
両者は $\mathcal U^*$ 内で一般に従属するので、評価像は $H-p$ と同型な部分線束となる。
商線束は $3H-B-(H-p)=2H-T$ であり、先と同じ右核の議論から
\[
                         d_C=H-T .
\]
特にこの値は、相手の二次曲線にも剰余点 $p$ にも依存しない。
\end{proof}

\begin{proof}[被覆と表示の既約性]
滑らかな行列式を持つ順序未指定の三平面対は
$\mathrm{Gr}(3,5)^2$ の非空な開部分で既約、次元12である。
楕円曲線の相対 $\mathrm{Pic}^0$ で $-2d=d_C$ を取ると有限エタール次数4の被覆となる。
上で求めた分裂 $G'$ と固定 $E$ の楕円五次曲線の開部分上でこの被覆は連結なので、
元の被覆も連結であり、滑らかな既約基底の上で既約である。
各点で $d$ の具体的な補助三行を行基底まで選ぶ自由度は $\mathrm{PGL}(U)$ の8である。
実際右核埋込みの座標を選ぶことで尽き、その全点を固定する射影変換はない。
さらに六次元線型空間の $\Gamma_Q$ が単射な非空開部分と、
交換子を混合行列に一致させる二つの順序付き基底の共通逆倍率の一変数を取る。
これらはエタール局所的に束の開部分ないしスカラー主束であり、
そこでの復元は線型方程式と定数階数の核による。従って全体は既約で次元 $12+8+6+1=27$ である。
非空な混合点の変形チャートはこの既約族に稠密である。
実際分裂点の構成で滑らかな三次式に出る方向があり、片側対称な正規化チャートも次元27である。
従って全ての持上げの所属成分と一般点の主張が定理\ref{thm:marked-mixed-component}から従う。

$E$ と計量を固定したとき、$G'$ の四基点は定理\ref{thm:marked-fixed-metric-elliptic}の
$\mathrm{Pic}^4$ 上の射影直線束のうち、二次鉛筆の完全交叉となる開部分で尽くされる。
逆にこの開部分では二次鉛筆の跡直交補が $G'$ を一意に復元し、
可換・非可換どちらでも相手三平面が実在する。非空性は例外外では一般の射影枠、
例外では平面を張る一般の三因子と三本の割線を避ける第四点で確かめられる。
この二次元基底上で、$F'$ の空間はGrassmann因子を五つの行列式の一次条件で切り、
行列式が非零の開部分を取ったものである。全繊維は既約一次元で、
全成分の次元は $2+6-5=3$ 以上である。よって縦成分はなく、全体は既約な次元3である
（一般の幾何繊維も前二定理の楕円曲線または二次曲線の開部分）。
例外外では一般の $G'$ が独立交換子となり、一般繊維の $-2$ 被覆も幾何学的既約である。
有限エタール被覆には縦成分が生じないため、全ての可換特殊値を含め同じ既約被覆に入る。
既知の混合変形の開チャートと同じ成分なので一般点は全次数正則である。
例外では全ての $T$ が同じ $M_0$ の切断であり、$d_C=H-M_0$ は固定された非零二等分点となる。
被覆はちょうどその四つの逆像ごとに分離する。どちらの場合も各表示成分の次元は $3+8+6+1=18$ である。
\end{proof}
